% Ad Familiares 15.12 --- headnote
% ad-familiares-15-12, September 51 BC, on the march through Lycaonia

\textit{Cicero --- already proconsul, en route through Lycaonia toward
his Cilician province in late September 51 BC --- to L. Aemilius Paullus,
who has just been elected consul for 50. The Perseus dateline groups
the letter with \textit{Fam.} 15.7, 15.8 and 15.9, all written on the
same stage of the eastward journey. The first half is a graceful note
of congratulation on Paullus's election, conventional in shape but
warm in tone: Cicero had never doubted the outcome, given Paullus's
services to the commonwealth and the standing of the Aemilian house,
but the news still produced ``incredible joy.''}

\textit{The second half is the practical request that runs through
almost every letter of this period to men of influence: Cicero wants
his term in Cilicia to be exactly one year and no longer. He had
accepted the province reluctantly, under the Pompeian law of 52, and
his terror is that he will be held there beyond his appointed year
through some neglect at Rome. He asks Paullus, on the strength of the
older man's zeal for him, to see that no injustice is done and that
nothing be added to his annual commission --- a refrain that will
become the keynote of his Cilician correspondence with the consuls of
50.}
